
\documentclass{article}
\usepackage{amsmath} % For advanced math features like \text, \frac, etc.
\usepackage[utf8]{inputenc} % To handle various characters
\begin{document}

\begin{table}[h]
\centering
\caption{Texas Temperatures Regression}
\label{tab:Texas-Temps}
\begin{tabular}{|c|c|c|c|c|} \hline 
Const & Lat & Elev & Long & Temp \\ \hline 
1.00000 &	29.7670 &	41.0000 &	95.3670 &	56.0000 \\ \hline 
1.00000 &	32.8500 &	440.000 &	96.8500 &	48.0000 \\ \hline 
1.00000 &	26.9330 &	25.0000 &	97.8000 &	60.0000 \\ \hline 
1.00000 &	31.9500 &	2851.00 &	102.183 &	46.0000 \\ \hline 
1.00000 &	34.8000 &	3840.00 &	102.467 &	38.0000 \\ \hline 
1.00000 &	33.4500 &	1461.00 &	99.6330 &	46.0000 \\ \hline 
1.00000 &	28.7000 &	815.000 &	100.483 &	53.0000 \\ \hline 
1.00000 &	32.4500 &	2380.00 &	100.533 &	46.0000 \\ \hline 
1.00000 &	31.8000 &	3918.00 &	106.400 &	44.0000 \\ \hline 
1.00000 &	34.8500 &	2040.00 &	100.217 &	41.0000 \\ \hline 
1.00000 &	30.8670 &	3000.00 &	102.900 &	47.0000 \\ \hline 
1.00000 &	36.3500 &	3693.00 &	102.083 &	36.0000 \\ \hline 
1.00000 &	30.3000 &	597.000 &	97.7000 &	52.0000 \\ \hline 
1.00000 &	26.9000 &	315.000 &	99.2830 &	60.0000 \\ \hline 
1.00000 &	28.4500 &	459.000 &	99.2170 &	56.0000 \\ \hline 
1.00000 &	25.9000 &	19.0000 &	97.4330 &	62.0000 \\ \hline 
\end{tabular}
\end{table}

\end{document}

